

\section{L'Input Compréhensible : Apprendre par le sens, pas par la forme}
space{0.3cm}
oindent

L'enseignement des langues anciennes, et singulièrement du grec ancien, se trouve aujourd'hui à la
croisée des chemins. Longtemps dominé par une tradition philologique axée sur l'analyse grammaticale
et la traduction (la méthode Grammaire-Traduction ou \textit{Grammar-Translation Method}, GTM), cet
enseignement fait face à une crise d'efficacité et de légitimité. Les étudiants, après plusieurs
années d'étude, peinent souvent à lire des textes originaux avec fluidité, restant dépendants du
dictionnaire et de l'analyse syntaxique consciente. Le point II.3.\cite{II-3-3-ref3} de votre projet
de refondation, qui postule que \og le cerveau n'apprend que par le sens (\textit{Meaning-focused
input}), pas par la forme (\textit{Form-focused instruction})\fg{}, ne constitue pas seulement une
prise de position méthodologique audacieuse ; il représente un alignement nécessaire de la pédagogie
des langues classiques sur plus de quarante années de recherche en acquisition des langues secondes
(\textit{Second Language Acquisition}, SLA), en psycholinguistique et en neurosciences cognitives.

Il s'agit de démontrer, preuves bibliographiques à l'appui, que l'exposition
à un \textit{Input Compréhensible} (CI) n'est pas une simple méthode parmi d'autres, mais le
mécanisme fondamental par lequel le cerveau humain traite et acquiert le langage. Nous explorerons
comment l'instruction focalisée sur la forme, caractéristique de l'approche traditionnelle, entre en
conflit avec l'architecture cognitive humaine, provoquant une surcharge de la mémoire de travail et
bloquant les processus de procéduralisation neuronale indispensables à la maîtrise réelle de la
langue.

L'analyse s'articulera autour de quatre axes majeurs. Premièrement, nous examinerons les théories
linguistiques de l'acquisition (Krashen, VanPatten) pour comprendre la primauté du sens dans le
traitement de l'information. Deuxièmement, nous plongerons dans les mécanismes neurobiologiques
(Modèle Déclaratif/Procédural d'Ullman) pour expliquer pourquoi l'enseignement explicite des règles
échoue à créer de la compétence implicite. Troisièmement, nous analyserons les contraintes
cognitives (Théorie de la Charge Cognitive de Sweller) et les seuils lexicaux (Nation) qui
conditionnent la lecture fluide. Enfin, nous évaluerons l'application concrète de ces principes à
travers les innovations pédagogiques contemporaines, notamment la méthode Polis et l'Expression
Séquentielle Vivante, tout en répondant aux critiques concernant la rigueur et la précision
linguistique.



\subsection{L'architecture cognitive de l'acquisition : la primauté du sens}

La pierre angulaire de la refondation pédagogique du grec ancien réside dans la compréhension des
mécanismes par lesquels le langage est internalisé par l'apprenant. La recherche en SLA a établi une
distinction fondamentale entre la connaissance consciente des règles d'une langue et la capacité
subconsciente à l'utiliser.



\subsection{L'hypothèse de l'input et la distinction acquisition/apprentissage}

Les travaux de Stephen Krashen dans les années 1980 ont posé les bases théoriques de l'approche
communicative moderne. Au cœur de son modèle se trouve la distinction, souvent débattue mais jamais
totalement réfutée dans ses implications pédagogiques, entre l'\textbf{acquisition} et
l'\textbf{apprentissage} (\textit{learning}).



\subsection{La nature subconsciente de l'acquisition}

L'acquisition est définie comme un processus subconscient, similaire, sinon identique, à la manière
dont les enfants développent leur langue maternelle (L1). L'apprenant n'est pas conscient du fait
qu'il acquiert une langue ; il est seulement conscient qu'il utilise la langue pour communiquer. Le
résultat de l'acquisition est une \og sensation\fg{} de correction grammaticale : les phrases \og
sonnent juste\fg{} ou non, sans que le locuteur puisse nécessairement invoquer la règle sous-
jacente.\cite{II-3-3-ref1} Selon Krashen, l'amélioration de la compétence linguistique dépend
\textit{uniquement} de l'acquisition et jamais de l'apprentissage conscient.\cite{II-3-3-ref3}



\subsection{Les limites de l'apprentissage conscient (le moniteur)}

À l'opposé, l'apprentissage est un processus conscient de développement de connaissances explicites
sur la langue : connaître les paradigmes de la troisième déclinaison grecque, les règles
d'accentuation, ou la syntaxe du génitif absolu. Krashen postule que ce savoir explicite a une
fonction très limitée : il agit uniquement comme un \og Moniteur\fg{} ou un éditeur. Il permet de
corriger la production après qu'elle a été initiée par le système acquis, et ce, uniquement si trois
conditions drastiques sont réunies : l'apprenant doit avoir du temps, il doit se focaliser sur la
forme, et il doit connaître la règle.\cite{II-3-3-ref2}

Dans le contexte de la lecture fluide du grec ancien, ces conditions sont rarement réunies.
L'analyse grammaticale consciente (l'usage du Moniteur) est cognitivement trop lente pour permettre
de suivre le fil d'un récit complexe comme l'Odyssée ou les Évangiles. En focalisant l'enseignement
sur la forme (Form-focused instruction), la pédagogie traditionnelle entraîne les élèves à devenir
des utilisateurs hyper-monitorés, capables de disséquer une phrase mot à mot mais incapables de la
lire.



\subsection{L'hypothèse de l'input (i+1)}

La conséquence pédagogique de cette distinction est l'Hypothèse de l'Input. Selon Krashen, nous
acquérons la langue d'une seule et unique manière : en comprenant des messages. L'acquisition se
produit lorsque l'apprenant reçoit un input compréhensible qui contient des structures légèrement
au-delà de son niveau de compétence actuel, noté i+1 (où i est le niveau actuel).\cite{II-3-3-ref1}

Pour le grec ancien, cela signifie que la confrontation avec des tableaux de conjugaison (input non
significatif) est inutile pour l'acquisition. Le cerveau doit être exposé à des phrases, des
histoires et des discours en grec dont le sens est rendu accessible par le contexte, les images ou
les gestes. Comme le soulignent Bailey et Fahad (2021) en revisitant Krashen, cette idée offre aux
enseignants un mandat clair : fournir aux apprenants des opportunités abondantes de \og créer du
sens\fg{} (making meaning) dans la langue cible.\cite{II-3-3-ref4}



\subsection{Le filtre affectif}

Krashen ajoute une dimension émotionnelle cruciale : l'Hypothèse du Filtre Affectif. La motivation,
la confiance en soi et l'anxiété jouent un rôle de modulateur. Un niveau élevé d'anxiété (fréquent
dans les cours de langues classiques où l'erreur est sanctionnée et la performance publique de
traduction exigée) \og lève\fg{} le filtre affectif, empêchant l'input, même compréhensible,
d'atteindre le dispositif d'acquisition du langage (\textit{Language Acquisition
Device}).\cite{II-3-3-ref1} L'approche communicative, en mettant l'accent sur le sens et en tolérant
les erreurs formelles initiales, tend à abaisser ce filtre, facilitant ainsi
l'acquisition.\cite{II-3-3-ref2}



\subsection{Le traitement de l'input (input processing) : la théorie de vanpatten}

Si Krashen a défini \textit{ce qu'il faut} pour apprendre (du sens), Bill VanPatten a explicité
\textit{comment} le cerveau traite cet input. Sa théorie du Traitement de l'Input (\textit{Input
Processing Theory}) est essentielle pour comprendre pourquoi l'instruction focalisée sur la forme
échoue souvent à produire des résultats.



\subsection{Le principe de la primauté du sens (primacy of meaning principle)}

VanPatten postule que les apprenants traitent l'input pour le sens \textit{avant} de le traiter pour
la forme. Ce principe repose sur la capacité limitée de la mémoire de travail humaine. Lorsqu'un
apprenant écoute ou lit une phrase en grec, sa priorité absolue est d'extraire le message
communicatif. S'il doit consacrer toutes ses ressources attentionnelles à comprendre les mots
lexicaux (le contenu sémantique), il ne dispose plus de ressources résiduelles pour traiter les
marqueurs grammaticaux (la forme).\cite{II-3-3-ref6}

Ce principe se décline en plusieurs sous-principes qui éclairent les difficultés spécifiques du grec
ancien :

\begin{table}[h!]\small\centering\begin{tabular}{| p{2.3cm} | p{3.5cm} | p{3.5cm} |}\hline\textbf{Sous-Principe} & \textbf{Description} & \textbf{Implication pour le Grec Ancien} \\ \hline\textbf{Primauté des Mots de Contenu} & L'apprenant cherche d'abord les noms et verbes porteurs de sens (ex: \og lion\fg{}, \og manger\fg{}).\cite{II-3-3-ref9} & Les particules, articles et désinences casuelles sont souvent ignorés au début car ils portent moins de charge sémantique immédiate. \\ \hline\textbf{Préférence Lexicale} & Si la grammaire et le lexique codent la même information, l'apprenant se fie au lexique.\cite{II-3-3-ref9} & Dans la phrase \og Hier (\textit{chthès}), il a marché (\textit{periepatès-e})\fg{}, le temps est marqué par l'adverbe et la désinence. L'apprenant traitera \og Hier\fg{} et ignorera la marque de l'aoriste, rendant l'enseignement explicite de la conjugaison inefficace tant que l'adverbe est présent. \\ \hline\textbf{Disponibilité des Ressources} & La forme n'est traitée que si la compréhension du sens ne sature pas l'attention.\cite{II-3-3-ref9} & Si le texte est trop difficile lexicalement, l'élève ne pourra jamais \og remarquer\fg{} (\textit{notice}) les structures grammaticales, bloquant l'acquisition syntaxique. \\ \hline\end{tabular}\end{table}

\subsection{Le principe du premier nom (first noun principle)}

Un autre obstacle majeur identifié par VanPatten est la tendance universelle des apprenants à
interpréter le premier nom ou pronom d'une phrase comme le sujet ou l'agent de
l'action.\cite{II-3-3-ref7}

Or, le grec ancien est une langue à ordre des mots libre, où l'objet (Accusatif) peut fréquemment
précéder le verbe ou le sujet. Si un élève lit la phrase Ton logon akouei ho maqètès (La parole \-
écoute \- le disciple), son cerveau, suivant le principe du Premier Nom, aura tendance à interpréter
\og La parole écoute...\fg{}, créant une confusion sémantique.

L'enseignement traditionnel tente de corriger cela par l'analyse logique (\og cherchez le
nominatif\fg{}). VanPatten suggère que cette approche lutte contre les mécanismes naturels de
traitement. La solution n'est pas l'explication grammaticale, mais l'Instruction de Traitement
(Processing Instruction) : structurer l'input de manière à ce que l'apprenant soit obligé de faire
attention à la forme (ici, la désinence de l'accusatif) pour comprendre le sens, sans pouvoir se
fier uniquement à l'ordre des mots ou au contexte.\cite{II-3-3-ref6}



\subsection{Focus on form vs. focus on forms : clarification conceptuelle}

Il est crucial de distinguer deux approches de l'enseignement de la forme, souvent confondues dans
les débats sur les langues classiques. Cette distinction, établie par Michael Long (1991), est
centrale pour valider le point II.3.\cite{II-3-3-ref3} de votre refondation.

1. \textbf{Focus on FormS (FonFS) :} C'est l'approche traditionnelle, typique de la méthode GTM. La
langue est découpée en unités discrètes (le génitif, l'optatif, la proposition infinitive) qui sont
enseignées séquentiellement comme des objets d'étude en soi. L'accent est mis primairement sur la
structure linguistique. Le \og syllabus\fg{} est une liste de points grammaticaux à
maîtriser.\cite{II-3-3-ref11}

2. \textbf{Focus on Form (FonF) :} C'est une approche où l'attention est attirée brièvement sur la
forme linguistique \textit{au cours d'une interaction dont l'objectif principal reste le sens}.
L'attention à la forme surgit incidemment, en réponse à un problème de communication ou de
compréhension, ou est planifiée de manière intégrée à une tâche communicative.\cite{II-3-3-ref11}

Les méta-analyses, notamment celles de Norris et Ortega (2000), ont montré que si l'instruction
explicite (FonFS) peut mener à des gains à court terme sur des tests de connaissances déclaratives,
elle est souvent moins efficace pour développer la compétence communicative fluide et
durable.\cite{II-3-3-ref12} De plus, le FonFS risque de créer des connaissances inertes : l'élève
\og sait\fg{} sa grammaire mais ne peut pas l'utiliser en temps réel.

Le point II.3.\cite{II-3-3-ref3} de votre projet, en rejetant l'instruction focalisée sur la forme
(Form-focused instruction au sens de FonFS), s'aligne sur les résultats montrant que l'acquisition
durable se produit lorsque l'attention à la forme est subordonnée à la quête du sens (Meaning-
focused).\cite{II-3-3-ref15}



\subsection{Les bases neurobiologiques : mémoire déclarative et procédurale}

L'argument le plus puissant en faveur de l'Input Compréhensible et contre l'approche grammaticale
traditionnelle ne vient pas de la pédagogie, mais des neurosciences. Le modèle Déclaratif/Procédural
(DP) développé par Michael Ullman fournit une explication biologique aux observations de Krashen et
VanPatten.



\subsection{La dissociation des systèmes de mémoire}

Ullman démontre que le langage chez l'humain repose sur deux systèmes de mémoire à long terme
distincts, qui ont des substrats neuronaux, des fonctions et des modes d'apprentissage
différents.\cite{II-3-3-ref17}



\subsection{La mémoire déclarative (le lexique)}

\begin{itemize}\item \textbf{Fonction :} Elle gère les connaissances explicites, les faits, et les événements (\og Savoir que\fg{}). Dans le langage, elle stocke le lexique (les mots, leurs sens), les expressions idiomatiques et les règles grammaticales apprises consciemment.\item \textbf{Substrat neuronal :} Elle est enracinée dans le lobe temporal médian, en particulier l'hippocampe.\cite{II-3-3-ref18}\item \textbf{Caractéristiques :} Elle permet un apprentissage très rapide (parfois en une seule exposition), mais l'accès à ces informations nécessite un contrôle attentionnel conscient. Elle est explicite et accessible à la conscience.\end{itemize}

\subsection{La mémoire procédurale (la grammaire mentale)}

\begin{itemize}\item \textbf{Fonction :} Elle gère les compétences motrices et cognitives automatisées (\og Savoir comment\fg{}). Dans le langage, elle sous-tend la grammaire mentale (syntaxe, morphologie, phonologie) et les séquences linguistiques combinatoires.\item \textbf{Substrat neuronal :} Elle repose sur les circuits fronto-striataux, incluant les ganglions de la base (noyau caudé) et l'aire de Broca dans le cortex frontal.\cite{II-3-3-ref20}\item \textbf{Caractéristiques :} Son apprentissage est lent et nécessite une répétition massive et une exposition constante à des stimuli (input). Elle fonctionne de manière implicite, rapide et automatique, sans effort conscient.\cite{II-3-3-ref21}\end{itemize}Pour un locuteur natif, la grammaire est procédurale. Il ne \og réfléchit\fg{} pas pour accorder un
adjectif ; son cerveau le fait automatiquement via les circuits striataux. En revanche,
l'enseignement traditionnel du grec ancien stocke la grammaire dans la \textbf{mémoire déclarative}.
L'élève apprend les règles comme des faits encyclopédiques. Pour lire, il doit récupérer ces faits
explicitement, ce qui est lent et cognitivement coûteux.\cite{II-3-3-ref19}



\subsection{L'effet de "bascule" (see-saw effect) et le blocage}

Le point crucial pour la refondation de l'enseignement est la découverte que ces deux systèmes ne
sont pas simplement indépendants, mais peuvent entrer en compétition. Ce phénomène est décrit par
Ullman comme l'effet de \og bascule\fg{} (\textit{see-saw effect}).\cite{II-3-3-ref23}



\subsection{Compétition et inhibition}

Des études suggèrent que l'activation intense de la mémoire déclarative (par exemple, lors de
l'apprentissage explicite de règles grammaticales complexes) peut inhiber ou supprimer l'activité de
la mémoire procédurale. L'œstrogène, par exemple, semble améliorer la mémoire déclarative tout en
déprimant la mémoire procédurale, illustrant les mécanismes biologiques de cette
balance.\cite{II-3-3-ref25}

Dans le contexte de l'apprentissage des langues (SLA), cela signifie que si l'on force un élève à
analyser consciemment la langue (approche Form-focused), on renforce le système déclaratif au
détriment du système procédural. On entraîne le cerveau à utiliser le \og mauvais\fg{} circuit pour
la grammaire.\cite{II-3-3-ref27}



\subsection{Le phénomène de blocage (\textit{blocking})}

Plus spécifiquement, l'instruction explicite peut provoquer un effet de \og blocage\fg{} (blocking)
sur l'apprentissage implicite. Si un apprenant possède une règle explicite pour une structure (par
exemple, la formation de l'aoriste), son cerveau peut cesser de prêter attention aux indices
statistiques subtils dans l'input qui permettraient normalement la procéduralisation de cette
structure.\cite{II-3-3-ref29} L'attention est détournée vers l'application de la règle consciente,
empêchant les mécanismes implicites de faire leur travail d'extraction de régularités.

Des études utilisant des potentiels évoqués (ERP) montrent que les apprenants formés explicitement
traitent les violations syntaxiques différemment des natifs (absence de la composante P600/LAN
caractéristique du traitement procédural), même après une pratique intensive.\cite{II-3-3-ref27} En
d'autres termes, l'enseignement explicite risque de fossiliser l'apprenant dans un mode de
traitement non-natif.



\subsection{La procéduralisation par l'input}

Le modèle DP prédit que pour atteindre une maîtrise native ou quasi-native (fluidité), la
connaissance grammaticale doit migrer de la mémoire déclarative vers la mémoire
procédurale.\cite{II-3-3-ref19} Cependant, cette \og procéduralisation\fg{} ne se fait pas par la
transformation directe d'une règle en automatisme, mais par le développement parallèle d'une
nouvelle compétence procédurale grâce à l'exposition massive à l'input et à la pratique en
contexte.\cite{II-3-3-ref31}

L'enseignement du grec doit donc viser à maximiser l'engagement des circuits procéduraux. Cela ne
peut se faire que par l'exposition à des séquences linguistiques en contexte (Meaning-focused
input), et non par l'étude de tableaux hors contexte.\cite{II-3-3-ref29} Le cerveau \og n'apprend
par le sens\fg{} que parce que le traitement du sens est la voie royale pour activer et entraîner
les ganglions de la base à prédire et traiter la structure grammaticale implicitement.



\subsection{La charge cognitive et les seuils lexicaux}

Si l'apprentissage passe par le sens, il est impératif que ce sens soit accessible sans épuiser les
ressources mentales de l'étudiant. La théorie de la Charge Cognitive (\textit{Cognitive Load
Theory}, CLT) de John Sweller offre un cadre analytique pour comprendre l'échec des méthodes
traditionnelles face aux textes complexes.



\subsection{La théorie de la charge cognitive (sweller)}

La mémoire de travail humaine est un goulot d'étranglement : elle ne peut traiter qu'un nombre très
limité d'éléments d'information nouveaux simultanément (environ 3 à 5 éléments, selon les études
récentes).\cite{II-3-3-ref37} La CLT distingue trois types de charge :

1. \textbf{Charge Intrinsèque :} Liée à la complexité inhérente du matériel (le grec ancien a une
charge intrinsèque élevée du fait de sa morphologie et de son alphabet).

2. \textbf{Charge Extrinsèque :} Liée à la manière dont l'information est présentée. Une mauvaise
pédagogie augmente cette charge inutilement.\cite{II-3-3-ref39}

3. \textbf{Charge Essentielle (Germane) :} La charge consacrée à la création de schémas mentaux et à
l'apprentissage.



\subsection{L'effet d'attention divisée (split-attention effect)}

Dans un cours de grec traditionnel, l'élève est souvent confronté à un texte original (charge
intrinsèque élevée) et doit constamment faire des allers-retours entre le texte, un dictionnaire, et
une grammaire. Sweller a identifié ce phénomène comme l'\og effet d'attention divisée\fg{} (split-
attention effect).\cite{II-3-3-ref41} L'élève doit intégrer mentalement des sources d'informations
disparates, ce qui sature sa mémoire de travail (charge extrinsèque massive).

Résultat : il ne reste aucune ressource cognitive pour la compréhension profonde du sens ou pour
l'acquisition incidente du vocabulaire. Le cerveau est en mode \og survie cognitive\fg{}, traitant
les mots un par un sans pouvoir construire une représentation globale du
discours.\cite{II-3-3-ref39}



\subsection{Le seuil critique de 98\% de couverture lexicale (hu \& nation)}

Pour éviter la surcharge cognitive et permettre un apprentissage focalisé sur le sens, il existe un
prérequis mathématique précis concernant le vocabulaire. Les recherches de Hu et Nation (2000),
confirmées par Schmitt (2011), ont établi que pour qu'une lecture soit fluide et permette la
compréhension sans assistance extérieure, le lecteur doit connaître \textbf{98\%} des mots du
texte.\cite{II-3-3-ref43}



\subsection{La dynamique des pourcentages}

\begin{itemize}\item \textbf{98\% de couverture :} (1 mot inconnu sur 50). La compréhension est optimale. Le contexte est suffisant pour deviner le sens du mot manquant (inférence). La charge cognitive est faible, permettant de se concentrer sur le message et le plaisir de lire (\textit{Reading for pleasure}). C'est la zone de l'Input Compréhensible idéal.\cite{II-3-3-ref46}\item \textbf{95\% de couverture :} (1 mot inconnu sur 20). La compréhension commence à se dégrader significativement pour la plupart des apprenants. L'inférence devient difficile.\cite{II-3-3-ref41}\item \textbf{90\% et moins :} (1 mot inconnu sur 10). C'est le seuil de frustration. Le lecteur ne peut plus suivre le fil de l'histoire. Il doit recourir à une stratégie de décodage laborieuse.\end{itemize}

\subsection{L'erreur des "textes authentiques"}

Le problème majeur de l'enseignement du grec est l'introduction prématurée de textes littéraires \og
authentiques\fg{} (Platon, Xénophon) qui présentent souvent une couverture lexicale bien inférieure
à 90\% pour un étudiant moyen. En forçant les élèves à travailler sur ces textes, on les place en
situation de surcharge cognitive permanente. Selon la CLT, cela bloque l'apprentissage car les
ressources sont entièrement consommées par la résolution de problèmes de bas niveau (identification
des mots), empêchant la construction de schémas en mémoire à long terme.\cite{II-3-3-ref49}

La refondation doit donc impérativement intégrer l'usage de textes gradués (graded readers) ou d'un
input oral contrôlé, garantissant ce seuil de 98\% pour permettre au cerveau d'apprendre par le
sens.



\subsection{Application pédagogique : vers une refondation du grec ancien}

Comment traduire ces principes théoriques (Primauté du Sens, Mémoire Procédurale, Charge Cognitive)
en une pratique pédagogique concrète pour le grec ancien? L'Institut Polis et d'autres chercheurs
ont développé des méthodologies qui répondent précisément à ce défi.



\subsection{La méthode polis et l'expression séquentielle vivante (lse)}

Sous la direction du professeur Christophe Rico, l'Institut Polis (Jérusalem) a élaboré une méthode
d'enseignement des langues anciennes comme des langues vivantes (\textit{Koine Greek as a living
language}), en s'appuyant directement sur les recherches en SLA.



\subsection{L'immersion et le tpr (total physical response)}

La méthode Polis utilise l'immersion totale dès le premier cours. Pour rendre l'input compréhensible
sans traduction (évitant ainsi l'effet d'attention divisée), elle utilise le \textit{Total Physical
Response} (TPR) développé par James Asher. L'enseignant donne des ordres en grec (\textit{\og
Anastèthi\fg{}} \- Lève-toi, \textit{\og Peripatès\fg{}} \- Marche) et les élèves exécutent
l'action.\cite{II-3-3-ref50}

\begin{itemize}\item \textbf{Justification cognitive :} Le TPR connecte directement le mot grec à l'action et à la sensation physique, créant une trace mnésique multimodale. Il contourne la mémoire déclarative (traduction) pour ancrer le sens dans une expérience vécue, favorisant une mémorisation rapide et durable du vocabulaire de base.\cite{II-3-3-ref53}\end{itemize}

\subsection{L'expression séquentielle vivante (living sequential expression \- lse)}

Innovation majeure de Rico, le LSE revisite les \og Séries\fg{} de François Gouin (XIXe siècle). Il
s'agit d'enseigner la langue à travers des séquences d'actions logiquement enchaînées qui racontent
une micro-histoire quotidienne.\cite{II-3-3-ref50}

\begin{itemize}\item \textbf{Exemple de Séquence (reconstituée) :}\end{itemize}1. \textit{Anistamai} (Je me lève)

2. \textit{Poreuomai pros tèn thuran} (Je marche vers la porte)

3. \textit{Anoigo tèn thuran} (J'ouvre la porte)

4. \textit{Exerchomai} (Je sors).\cite{II-3-3-ref55}

\begin{itemize}\item \textbf{Avantages Neurocognitifs du LSE :}\item \textbf{Réduction de la Charge Cognitive :} La logique causale et chronologique de la séquence aide la mémoire de travail à prédire l'étape suivante. L'élève n'a pas à mémoriser une liste arbitraire de mots, mais un \og script\fg{} d'action.\cite{II-3-3-ref57}\item \textbf{Contextualisation Grammaticale (Form-Meaning Mapping) :} La grammaire est apprise par la variation de la séquence. L'enseignant demande : \og Que fait-il?\fg{} (Troisième personne). L'élève répond : \textit{\og Anistatai\fg{}} (Il se lève). La variation morphologique (-mai / \-tai) est associée directement au changement d'acteur, sans passer par un tableau abstrait. C'est l'application parfaite du \textit{Focus on Form} (attention à la forme en contexte) préconisé par Long.\cite{II-3-3-ref54}\item \textbf{Facilitation de la Parole :} En permettant aux étudiants de verbaliser des actions concrètes, le LSE active les aires motrices du langage (Broca), renforçant les circuits procéduraux nécessaires à la fluidité.\cite{II-3-3-ref58}\end{itemize}

\subsection{L'approche des "quatre fils" (the four strands) de paul nation}

Pour structurer un curriculum complet qui respecte l'équilibre entre sens et forme, la refondation
peut s'appuyer sur le cadre des \og Quatre Fils\fg{} de Paul Nation. Nation soutient qu'un cours de
langue équilibré doit consacrer du temps égal (environ 25\% chacun) à quatre types d'activités 59 :

\begin{table}[h!]\small\centering\begin{tabular}{| p{2.3cm} | p{3.5cm} | p{3.5cm} |}\hline\textbf{Le Fil (Strand)} & \textbf{Description \& Application au Grec Ancien} & \textbf{Objectif Cognitif} \\ \hline\textbf{Meaning-focused Input} (Input axé sur le sens) & Écoute et lecture de textes accessibles (98\% connus). Utilisation de \textit{graded readers} (textes simplifiés). & Acquisition incidente de vocabulaire et de grammaire. Consolidation des schémas. \\ \hline\textbf{Meaning-focused Output} (Output axé sur le sens) & Parler et écrire pour communiquer un message réel. Utiliser le grec pour décrire des images, raconter des histoires (méthode Polis). & Renforcement des connexions neuronales productives. Le \og Pushed Output\fg{} (Swain) force à traiter la syntaxe plus profondément.\cite{II-3-3-ref62} \\ \hline\textbf{Language-focused Learning} (Apprentissage de la langue) & Étude délibérée du vocabulaire (flashcards), explications grammaticales ponctuelles (FonF). & Connaissance explicite qui peut aider à \og remarquer\fg{} (\textit{noticing}) les formes dans l'input futur. Ne doit pas dépasser 25\% du temps. \\ \hline\textbf{Fluency Development} (Développement de la fluidité) & Activités rapides sur du matériel \textit{très facile} (100\% connu). Relecture de textes déjà étudiés, écoute répétée. & Automatisation des connaissances (procéduralisation). Augmentation de la vitesse d'accès lexical.\cite{II-3-3-ref64} \\ \hline\end{tabular}\end{table}L'erreur tragique de l'enseignement traditionnel est de consacrer près de 100\% du temps au
\textit{Language-focused Learning} (analyse et traduction), négligeant totalement les trois autres
fils, en particulier l'Input axé sur le sens et le Développement de la fluidité.



\subsection{Réponses aux critiques et débats}

L'adoption de ces méthodes communicatives pour une langue dite \og morte\fg{} et au corpus clos
suscite des réserves légitimes qu'il convient d'adresser scientifiquement.



\subsection{La critique du "pidgin" et de l'imprécision}

Une critique récurrente est que l'enseignement communicatif du grec produirait un \og pidgin\fg{} ou
une langue approximative, sacrifiant la précision morphologique à la fluidité superficielle. Les
détracteurs craignent que sans l'étude rigoureuse des tableaux, les élèves ne maîtrisent jamais la
complexité des déclinaisons.\cite{II-3-3-ref66}

\begin{itemize}\item \textbf{Réponse Scientifique :} Il est vrai que l'immersion \textit{sans} instruction formelle peut mener à une fossilisation des erreurs. C'est pourquoi la recherche actuelle privilégie le \textit{Focus on Form} (et non le pur \textit{Meaning-focused} sans aucune correction). Cependant, les études montrent que les élèves issus de méthodes communicatives bien structurées finissent par atteindre une précision égale ou supérieure, car leur savoir est opératoire et non seulement théorique.\cite{II-3-3-ref69}\item \textbf{L'Argument du \og Pidgin\fg{} vs. GTM :} Il faut comparer ce qui est comparable. La méthode GTM produit-elle de la précision? Elle produit une capacité à \textit{analyser} la forme, mais souvent une incapacité totale à \textit{produire} une phrase correcte spontanément. Un \og interlangue\fg{} temporairement imparfait (comme le pidgin d'un apprenant) est une étape naturelle et nécessaire vers la maîtrise, prévue par l'Hypothèse de l'Ordre Naturel de Krashen. Vouloir une précision parfaite dès le début est un frein cognitif qui bloque la fluidité (Filtre Affectif élevé).\cite{II-3-3-ref70}\end{itemize}

\subsection{La complexité morphologique spécifique du grec}

Le grec ancien, avec sa morphologie flexionnelle riche, est-il inadapté à l'approche communicative
conçue pour l'anglais ou l'espagnol?

\begin{itemize}\item \textbf{Réponse :} Au contraire. La richesse morphologique impose une charge cognitive très lourde à la mémoire déclarative. Il est impossible de \og penser\fg{} consciemment à toutes les désinences en temps réel. Seule la \textbf{mémoire procédurale}, capable de gérer des traitements statistiques complexes implicitement, peut maîtriser une telle morphologie. Ullman souligne que plus la structure est complexe, plus elle bénéficie de la procéduralisation par l'input massif.\cite{II-3-3-ref19} Le LSE de Rico, en associant geste et morphologie, est particulièrement adapté pour \og faire sentir\fg{} la valeur des cas (accusatif directionnel vs locatif) sans passer par l'abstraction.\cite{II-3-3-ref53}\end{itemize}

\subsection{Le statut de "langue morte"}

Peut-on traiter le grec ancien comme une langue vivante alors qu'il n'y a pas de locuteurs natifs?

\begin{itemize}\item \textbf{Réponse Neurocognitive :} Pour le cerveau, la distinction \og langue morte\fg{} / \og langue vivante\fg{} n'existe pas. Il n'y a que des stimuli linguistiques. Si l'input est riche, interactif et signifiant, le cerveau active les mêmes mécanismes d'acquisition (aires de Broca, ganglions de la base) que pour n'importe quelle autre langue.\cite{II-3-3-ref55} Créer une communauté de locuteurs en classe (comme le fait l'Institut Polis) réactive le potentiel \og vivant\fg{} de la langue dans l'esprit de l'apprenant.\end{itemize}

\subsection{Conclusion et synthèse des tableaux}

La recherche approfondie valide pleinement le point II.3.\cite{II-3-3-ref3} de votre projet de
refondation. L'affirmation selon laquelle \og le cerveau n'apprend que par le sens\fg{} est soutenue
par la convergence des modèles théoriques (Krashen, VanPatten), des preuves neurobiologiques
(Ullman) et des contraintes cognitives (Sweller).

L'obstination pédagogique sur la forme (\textit{Form-focused instruction}) constitue une aberration
cognitive qui sature la mémoire de travail, induit un traitement déclaratif lent et bloque le
développement des automatismes procéduraux nécessaires à la lecture fluide. À l'inverse, une
pédagogie refondée sur l'\textit{Input Compréhensible}, structurée par des méthodes comme le LSE et
respectant les seuils de couverture lexicale, offre la seule voie scientifiquement cohérente pour
former des lecteurs compétents de grec ancien.



\subsection{Tableaux récapitulatifs}



\subsection{Tableau 1 : impact neurocognitif comparé des méthodes pédagogiques}

\begin{table}[h!]\small\centering\begin{tabular}{| p{2cm} | p{2.2cm} | p{2.2cm} | p{2.2cm} |}\hline\textbf{Dimension} & \textbf{Méthode Traditionnelle (GTM)} & \textbf{Refondation (Input Compréhensible / Polis)} & \textbf{Justification Scientifique} \\ \hline\textbf{Priorité Cognitive} & Forme (Règles, Paradigmes) & Sens (Message, Action) & \textit{Primacy of Meaning} (VanPatten) : Le cerveau traite le sens avant la forme. \\ \hline\textbf{Système Mémoire} & Déclaratif (Hippocampe) & Procédural (Ganglions de la Base) & Modèle DP (Ullman) : La fluidité dépend de la mémoire procédurale. \\ \hline\textbf{Charge Cognitive} & Surcharge Extrinsèque (Attention divisée) & Optimisée (Charge intrinsèque gérée) & CLT (Sweller) : L'attention divisée bloque l'apprentissage. \\ \hline\textbf{Traitement} & Traduction / Décodage (Lent) & Lecture / Compréhension directe (Rapide) & La traduction est une tâche de résolution de problème, pas de lecture (Psycholinguistique). \\ \hline\textbf{Résultat neuronal} & Traitement non-natif (N400 seulement) & Traitement quasi-natif (P600/LAN) & Les ERP montrent que l'instruction explicite ne crée pas les signatures neuronales natives.\cite{II-3-3-ref27} \\ \hline\end{tabular}\end{table}

\subsection{Tableau 2 : les piliers de la refondation selon les données de la recherche}

\begin{table}[h!]\small\centering\begin{tabular}{| p{2.3cm} | p{3.5cm} | p{3.5cm} |}\hline\textbf{Pilier Pédagogique} & \textbf{Application Concrète} & \textbf{Fondement Recherche} \\ \hline\textbf{Input Massif \& Signifiant} & 75\% du temps de cours dédié à l'écoute/lecture de textes compris à 98\%. & Hypothèse de l'Input (Krashen), Seuil de 98\% (Hu \& Nation). \\ \hline\textbf{Immersion \& Action} & TPR et LSE (Séquences d'actions verbalisées). Usage exclusif du grec. & Réduction du filtre affectif, Ancrage sensori-moteur (Rico/Asher). \\ \hline\textbf{Grammaire Contextualisée} & Focus on Form (FonF) : explication brève \textit{après} ou \textit{pendant} l'usage. & Éviter le \og blocage\fg{} de l'apprentissage implicite (Ellis/Ullman). \\ \hline\textbf{Fluidité} & Lectures répétées de textes faciles (\textit{Easy reading}). & Automatisation procédurale (DeKeyser/Nation). \\ \hline\end{tabular}\end{table}La refondation de l'enseignement du grec ancien n'est donc pas une simple modernisation de surface,
mais une nécessité scientifique pour aligner la didactique sur la réalité biologique de
l'apprentissage humain.

